
\documentclass{article}
\usepackage{colortbl}
\usepackage{makecell}
\usepackage{multirow}
\usepackage{supertabular}

\begin{document}

\newcounter{utterance}

\centering \large Interaction Transcript for game `cladder', experiment `full\_v1.5\_default', episode 6916 with reflexion{-}llama{-}3.1{-}8b{-}instruct.
\vspace{24pt}

{ \footnotesize  \setcounter{utterance}{1}
\setlength{\tabcolsep}{0pt}
\begin{supertabular}{c@{$\;$}|p{.15\linewidth}@{}p{.15\linewidth}p{.15\linewidth}p{.15\linewidth}p{.15\linewidth}p{.15\linewidth}}
    \# & \multicolumn{2}{c}{Player} && \multicolumn{2}{c}{Game Master} \\
    \hline

    \theutterance \stepcounter{utterance}  
    & & & \multicolumn{4}{p{0.6\linewidth}}{
        \cellcolor[rgb]{0.9,0.9,0.9}{
            \makecell[{{p{\linewidth}}}]{
                \texttt{\tiny{[P1$\langle$GM]}}
                \texttt{You are an expert in causal inference. The following question is not a typical commonsense query, but rather a meticulously designed question created by a professor specializing in causal inference, intended to assess the students' mastery of the course content.} \\
\texttt{Imagine a self{-}contained, hypothetical world with only the following conditions, and without any unmentioned factors or causal relationships: Jyka has a direct effect on hwax. Lirg has a direct effect on hwax. Method 1: We look at how jyka correlates with lirg case by case according to hwax. Method 2: We look directly at how jyka correlates with lirg in general. To understand how jyka affects lirg, is it more correct to use the Method 1 than Method 2?} \\
\texttt{Start your answer with 'Yes' or 'No', followed by additional reasoning or evidence to support your explanation.} \\
            }
        }
    }
     \\ \\

    \theutterance \stepcounter{utterance}  
    & \multicolumn{4}{p{0.6\linewidth}}{
        \cellcolor[rgb]{0.561,0.737,0.561}{
            \makecell[{{p{\linewidth}}}]{
                \texttt{\tiny{[P1$\rangle$GM]}}
                \texttt{No, it is not more correct to use Method 1 than Method 2 to understand how jyka affects lirg.} \\
\\ 
\texttt{In this hypothetical world, we have two direct causal relationships: Jyka has a direct effect on hwax, and Lirg has a direct effect on hwax. This creates a situation where the effect of Jyka on Hwax may be confounded by the effect of Lirg on Hwax. } \\
\\ 
\texttt{When we look at the correlation between Jyka and Lirg case by case according to Hwax (Method 1), we are essentially controlling for the effect of Hwax on both Jyka and Lirg. However, this control may not be sufficient to eliminate the confounding effect of Lirg on Hwax, which in turn affects Jyka. This is because the relationship between Jyka and Lirg may be influenced by the level of Hwax, even after controlling for it.} \\
\\ 
\texttt{On the other hand, when we look directly at the correlation between Jyka and Lirg in general (Method 2), we are not controlling for any intermediate variables like Hwax. This approach is more likely to capture the true causal relationship between Jyka and Lirg, as it does not introduce any additional assumptions or biases.} \\
\\ 
\texttt{In general, when dealing with confounding variables, it is often better to use a more direct approach, such as Method 2, rather than trying to control for intermediate variables, as in Method 1. This is because controlling for intermediate variables can sometimes introduce new biases or assumptions, rather than eliminating them.} \\
            }
        }
    }
    & & \\ \\

    \theutterance \stepcounter{utterance}  
    & & & \multicolumn{2}{p{0.3\linewidth}}{
        \cellcolor[rgb]{0.9,0.9,0.9}{
            \makecell[{{p{\linewidth}}}]{
                \texttt{\tiny{[GM$|$GM]}}
                \texttt{No} \\
            }
        }
    }
    & & \\ \\

    \theutterance \stepcounter{utterance}  
    & & & \multicolumn{2}{p{0.3\linewidth}}{
        \cellcolor[rgb]{0.9,0.9,0.9}{
            \makecell[{{p{\linewidth}}}]{
                \texttt{\tiny{[GM$|$GM]}}
                \texttt{no} \\
            }
        }
    }
    & & \\ \\

    \theutterance \stepcounter{utterance}  
    & & & \multicolumn{2}{p{0.3\linewidth}}{
        \cellcolor[rgb]{0.9,0.9,0.9}{
            \makecell[{{p{\linewidth}}}]{
                \texttt{\tiny{[GM$|$GM]}}
                \texttt{game\_result = WIN} \\
            }
        }
    }
    & & \\ \\

\end{supertabular}
}

\end{document}
